\documentclass[11pt]{article}
\usepackage[a4paper,margin=25mm]{geometry}
\usepackage[T1]{fontenc}
\usepackage[utf8]{inputenc}
\usepackage{lmodern}
\usepackage{amsmath,amssymb,amsthm}
\usepackage{hyperref}
\hypersetup{hidelinks}
\usepackage{listings}
\lstset{basicstyle=\ttfamily\small,breaklines=true,columns=fullflexible,keepspaces=true,
 literate={ℝ}{{$\mathbb R$}}1 {ℕ}{{$\mathbb N$}}1 {ℤ}{{$\mathbb Z$}}1
 {∧}{{$\land$}}1 {∃}{{$\exists$}}1 {≠}{{$\ne$}}1 {∀}{{$\forall$}}1
 {→}{{$\to$}}1 {≤}{{$\le$}}1 {∈}{{$\in$}}1 {¬}{{$\neg$}}1
 {⁻¹}{{$^{-1}$}}2}
\newtheorem{theorem}{Theorem}
\newtheorem{lemma}[theorem]{Lemma}
\newtheorem{corollary}[theorem]{Corollary}
\newcommand{\QI}{\operatorname{QI}}
\newcommand{\HH}{\operatorname{Hyp}}
\title{Disproof of TLMC Conjecture 00000007713:\\Incompatible Quantitative Clauses for Corona Growth}
\author{AlyciaBHZ on behalf of the Omega Institute}
\date{5 October 2026}
\begin{document}
\maketitle

\begin{abstract}
The conjecture requires an exponential growth rate $\lambda(p,q)$ to satisfy
$\lambda+\lambda^{-1}=(p-2)(q-2)-2$, while allowing quadratic irrationality
only if $(p-2)(q-2)\in\{4,5,6,10\}$. These exact clauses are incompatible.
At the hyperbolic pair $(4,6)$ the equation has the two quadratic irrational
roots $3\pm2\sqrt2$, although $(4-2)(6-2)=8$ is excluded by the list.
More generally, for every integer $N\ge3$, every real solution of
$x+x^{-1}=N$ is quadratic irrational. Consequently the same contradiction
occurs on every simultaneous tail $p,q\ge K$, using $p=q=\max(K,6)$.
The global and tail refutations are both proved in Lean 4 with Mathlib;
no knowledge of the true corona growth rate is needed.
\end{abstract}

\section{Statement and interpretation}
The English statement of \textbf{Conjecture 00000007713} is reproduced
verbatim below (mathematical symbols are typeset):
\begin{quote}
Definition: For the regular tiling $\{p,q\}$ of the hyperbolic plane, the vertex
corona count $C_k$ is the number of vertices at graph distance at most $k$
from a fixed vertex. Conjecture: For the two-parameter family with $p,q$ both
tending to infinity, the growth rate of $C_k$ is given by a quadratic form
in $2\cosh$, with exponential growth rate $\lambda(p,q)$ satisfying the
transcendental equation $\lambda+\lambda^{-1}=(p-2)(q-2)-2$, and $\lambda$ is
a quadratic irrational only when $(p-2)(q-2)$ lies in $\{4,5,6,10\}$.
(transcendental corona growth equation)
\end{quote}

The regular hyperbolic tiling parameters are integers $p,q\ge3$ with
\[
\HH(p,q)\quad\Longleftrightarrow\quad
p\ge3,\quad q\ge3,\quad \frac1p+\frac1q<\frac12.
\]
Write $S(p,q)=(p-2)(q-2)$ and $L=\{4,5,6,10\}$. We use the standard definition
\[
\QI(x)\quad\Longleftrightarrow\quad
x\notin\mathbb Q\ \text{ and }\
\exists a,b,c\in\mathbb Z:\ a\ne0,\quad ax^2+bx+c=0.
\]
An irrational root of a genuine quadratic over $\mathbb Q$ has algebraic
degree exactly two; this definition therefore has the usual mathematical
meaning of ``quadratic irrational.''

The two precise quantitative clauses in the statement are
\begin{align}
\tag{A}\label{A}
&\forall p,q\text{ with }\HH(p,q),\qquad
\lambda(p,q)+\lambda(p,q)^{-1}=S(p,q)-2,\\
\tag{B}\label{B}
&\forall p,q\text{ with }\HH(p,q),\qquad
\QI(\lambda(p,q))\ \Longrightarrow\ S(p,q)\in L.
\end{align}
``Only when'' expresses the implication in (B), not its converse.
The written equality in (A) is read literally. The phrase ``quadratic form
in $2\cosh$'' lacks a specified form, but any further condition on the growth
rate cannot make these two necessary clauses compatible. This argument
neither defines a replacement corona growth rate nor assumes that the
claimed equation correctly describes it.

The reference to $p,q$ tending to infinity admits a second natural reading:
(A) and (B) are asserted only once both parameters are sufficiently large.
We refute that reading too: for every cutoff $K$, no real-valued function
satisfies both clauses on all hyperbolic pairs $p,q\ge K$. In fact, on the
diagonal every individual pair $(t,t)$ with $t\ge6$ gives the contradiction.
Thus neither a finite exceptional range nor restriction to the large diagonal
repairs these exact clauses. An equality asserted only approximately or in a
limit would be a different statement, requiring an error term or a specified
limit; such a modified statement is not asserted or refuted here.

\section{The elementary arithmetic obstruction}
\begin{lemma}\label{nonsquare}
For every integer $N\ge3$, the integer $N^2-4$ is not a square.
\end{lemma}
\begin{proof}
We have
\[
(N-1)^2<N^2-4<N^2,
\]
because the difference in the left inequality is $2N-5>0$ and the difference
in the right inequality is $4>0$. If $m^2=N^2-4$ for an integer $m$, taking
absolute values would give $N-1<|m|<N$. No integer lies strictly between the
consecutive integers $N-1$ and $N$.
\end{proof}

\begin{lemma}\label{rational-square}
If a rational number $r$ has integer square, then $r$ is an integer.
\end{lemma}
\begin{proof}
Write $r=a/b$ with $a\in\mathbb Z$, $b\in\mathbb Z_{>0}$, and
$\gcd(|a|,b)=1$. If $r^2=d\in\mathbb Z$, then $a^2=db^2$. If a prime
$\ell$ divided $b$, it would divide $a^2$ and hence $a$, contradicting
coprimality. Thus $b$ has no prime divisor, so $b=1$.
\end{proof}

\begin{theorem}\label{general}
Let $N\in\mathbb Z$ satisfy $N\ge3$, and let $x\in\mathbb R$ satisfy
$x+x^{-1}=N$. Then $x$ is a quadratic irrational.
\end{theorem}
\begin{proof}
For the ordinary reciprocal, $x\ne0$ is implicit. Lean defines $0^{-1}=0$;
under this convention the equation itself excludes $x=0$, since $N\ge3$.
Multiplying the equation by $x$ gives
\begin{equation}\label{quadratic}
x^2-Nx+1=0.
\end{equation}
This is an integer quadratic with nonzero leading coefficient $1$.
Completing the square gives
\[
(2x-N)^2=N^2-4.
\]
If $x$ were rational, $2x-N$ would be rational with integer square.
By Lemma~\ref{rational-square} it would be an integer, contradicting
Lemma~\ref{nonsquare}. Hence $x$ is irrational and satisfies
\eqref{quadratic}, proving $\QI(x)$.

Equivalently, the two roots are
\[
x=\frac{N\pm\sqrt{N^2-4}}{2}.
\]
Both are positive, reciprocal, and irrational; the plus root is greater
than $1$ and the minus root is less than $1$. The argument proves
irrationality for both without selecting the growth branch.
\end{proof}

For context, every hyperbolic pair in the specified integer range has
$S(p,q)\ge5$. Indeed, multiplying $1/p+1/q<1/2$ by the positive number $2pq$
gives $2(p+q)<pq$, and therefore
\[
S(p,q)=pq-2(p+q)+4>4.
\]
Since $S(p,q)$ is an integer, $N=S(p,q)-2\ge3$. Thus if the conjectured
equation holds, its solution is quadratic irrational at \emph{every}
hyperbolic pair. The formal proof below needs only the explicit pair
$(4,6)$ and the diagonal $t\ge6$; this additional observation explains why
the arithmetic restriction is systematically incompatible with the equation.

\section{Concrete and arbitrarily large contradictions}
\begin{theorem}\label{concrete}
There is no function $\lambda:\mathbb N\times\mathbb N\to\mathbb R$
satisfying both (A) and (B).
\end{theorem}
\begin{proof}
The tiling parameters $(p,q)=(4,6)$ are hyperbolic:
\[
\frac14+\frac16=\frac5{12}<\frac12.
\]
Here $S=(4-2)(6-2)=8$. Clause (A) yields
\[
\lambda(4,6)+\lambda(4,6)^{-1}=6.
\]
By Theorem~\ref{general}, $\lambda(4,6)$ is quadratic irrational.
In this case the quadratic is $x^2-6x+1=0$, with roots
$3\pm2\sqrt2$. Clause (B) would then force
$8\in\{4,5,6,10\}$, which is false.
\end{proof}

\begin{theorem}\label{tail}
For every $K\in\mathbb N$ and every function
$\lambda:\mathbb N\times\mathbb N\to\mathbb R$, the restrictions of
(A) and (B) to hyperbolic pairs with $p,q\ge K$ cannot both hold.
\end{theorem}
\begin{proof}
Let $t=\max(K,6)$ and take $p=q=t$. Both parameters are at least $K$,
and the pair is hyperbolic since
\[
\frac1t+\frac1t=\frac2t\le\frac13<\frac12.
\]
Let $M=(t-2)^2$. Since $t\ge6$, we have $M\ge16$, so $M\notin L$.
Set $N=M-2\ge14\ge3$. The restricted clause (A) gives
$\lambda(t,t)+\lambda(t,t)^{-1}=N$, whence
$\QI(\lambda(t,t))$ by Theorem~\ref{general}.
The restricted clause (B) forces $M\in L$, a contradiction.
\end{proof}

In particular there is no existentially chosen cutoff $K$ beyond which the
clauses are simultaneously valid. If the two clauses are assigned separate
cutoffs, their maximum is a common cutoff, so that reading fails as well.
The contradiction also applies when the growth rate is required to be
positive or greater than $1$, since the theorems already quantify over all
real-valued functions and prove impossibility without additional hypotheses.
A stronger reading of ``only when'' as an equivalence still includes (B)
and is therefore refuted.

\section{Lean formalization and semantic correspondence}
The project is pinned to \texttt{leanprover/lean4:v4.33.0} and Mathlib tag
\texttt{v4.33.0}, whose manifest revision is
\texttt{db584cd6d46c92f209a44c0f1c829460d327499d}.
The sole source module is \texttt{lean4/CoronaGrowth.lean}. The following
statements reproduce the definitions and main theorem from the source:
\begin{lstlisting}
def IsQuadraticIrrational (x : ℝ) : Prop :=
  Irrational x ∧ ∃ a b c : ℤ, a ≠ 0 ∧
    (a : ℝ) * x ^ 2 + (b : ℝ) * x + (c : ℝ) = 0

def Hyperbolic (p q : ℕ) : Prop :=
  3 ≤ p ∧ 3 ≤ q ∧ 1 / (p : ℝ) + 1 / (q : ℝ) < 1 / 2

def EquationClause (lam : ℕ → ℕ → ℝ) : Prop :=
  ∀ p q, Hyperbolic p q →
    lam p q + (lam p q)⁻¹ = ((p : ℝ) - 2) * ((q : ℝ) - 2) - 2

def OnlyWhenClause (lam : ℕ → ℕ → ℝ) : Prop :=
  ∀ p q, Hyperbolic p q → IsQuadraticIrrational (lam p q) →
    (p - 2) * (q - 2) ∈ ({4, 5, 6, 10} : Finset ℕ)

theorem conjecture_00000007713_false (lam : ℕ → ℕ → ℝ) :
    ¬ (EquationClause lam ∧ OnlyWhenClause lam)
\end{lstlisting}
The definitions in the main theorem expand to exactly (A) and (B).
Natural subtraction in the membership clause agrees with integer subtraction
because every \texttt{Hyperbolic} pair has $p,q\ge3$; it does not truncate any
quantity used in the proof. The equation uses real subtraction explicitly.

The more robust theorem states:
\begin{lstlisting}
theorem conjecture_00000007713_false_on_every_tail
    (K : ℕ) (lam : ℕ → ℕ → ℝ) : ¬ TailClauses K lam

theorem no_eventual_quantitative_clauses (lam : ℕ → ℕ → ℝ) :
    ¬ ∃ K : ℕ, TailClauses K lam
\end{lstlisting}
Here \texttt{TailClauses K lam} is precisely the conjunction of (A) and (B),
with \texttt{K <= p} and \texttt{K <= q} hypotheses added to each clause.
Thus the tail theorem formalizes Theorem~\ref{tail}, and the final theorem
rules out choosing any sufficient cutoff.

The supporting theorem is:
\begin{lstlisting}
theorem equation_quadratic_irrational (N : ℤ) (hN : 3 ≤ N) (x : ℝ)
    (heq : x + x⁻¹ = (N : ℝ)) : IsQuadraticIrrational x
\end{lstlisting}
Its proof follows Theorem~\ref{general}. The theorem
\path{discriminant_not_square} proves Lemma~\ref{nonsquare} over the actual
integers. The imported standard Mathlib theorem
\texttt{Rat.isSquare\_intCast\_iff} supplies the rational-square fact of
Lemma~\ref{rational-square}; the proof uses rational and real casts only to
transfer equalities, not as approximations. The quadratic coefficients
witnessed in Lean are $a=1$, $b=-N$, and $c=1$.
\texttt{equation\_nonzero} handles Lean's total inverse explicitly.
The two counterexample proofs invoke these lemmas and certify all finite
arithmetic and membership facts within Lean.

No graph or tiling construction is required in Lean: the claim assigns a
real growth rate to every eligible parameter pair and entails the two
formalized numerical conditions. The Lean theorem proves that no function
at all can meet those conditions on the specified domain. Therefore any
actual function of corona growth rates also cannot meet them. This is a
refutation by internal inconsistency, rather than a calculation of the true
corona counts or a weakening of either exact quantitative clause.

\section{Verification and axiom audit}
The final source compiled successfully with Lean and Mathlib \texttt{v4.33.0}.
The compiler output and axiom audit are saved in \texttt{verification.txt}.
The source ends with \texttt{\#print axioms} commands auditing the general
arithmetic lemmas and all three refutation theorems. Every audited theorem
has the same axiom output:
\[
\texttt{[propext, Classical.choice, Quot.sound]}.
\]
These are the permitted standard Lean axioms. There are no added axioms,
no incomplete proofs, and no use of \texttt{native\_decide}. Arithmetic tactics
construct proofs checked by the Lean kernel.

The optional standard-library script \texttt{verify.py} uses exact rational
arithmetic and integer square roots. It checks the $(4,6)$ instance and both
roots in $\mathbb Q(\sqrt2)$; the discriminant inequalities for
$3\le N\le10000$; all hyperbolic pairs $3\le p,q\le100$; and the diagonal
witnesses for $0\le K\le1000$. Its output is also included in
\texttt{verification.txt}. These finite computations are independent sanity
checks. The infinite theorems are established by the mathematical proof
and Lean, with no dependence on the script.

To build the delivered project in a fresh environment with Lean installed:
\begin{lstlisting}
cd lean4
lake exe cache get && lake build
\end{lstlisting}
To rerun the optional checks and rebuild the report from the package root:
\begin{lstlisting}
python3 verify.py
pdflatex -interaction=nonstopmode -halt-on-error report.tex
pdflatex -interaction=nonstopmode -halt-on-error report.tex
\end{lstlisting}

\section{Scope and reading choices}
The disproof covers exact equality, the necessary implication meant by
``only when,'' all valid hyperbolic parameters, every simultaneous tail,
and every sufficiently large diagonal pair. It requires no additional
hypothesis on the value of $\lambda$. The assertion about quadratic forms
in $2\cosh$ and the actual behavior of $C_k$ remain unspecified, but cannot
repair the incompatible necessary clauses. The observation that all
hyperbolic pairs have $N\ge3$ is proved in this report; Lean formalizes the
general integer-$N$ lemma and all parameter facts needed for both main
refutations, without formalizing that additional all-pairs observation.

A reading that replaces the displayed equality by an asymptotic approximation
is not a literal reading of the exact equation. This submission makes no
claim about such an unstated replacement. The large-parameter theorem
addresses the explicit ``both tending to infinity'' qualification while
preserving the two exact quantitative assertions.

\medskip
\noindent\textbf{Attribution.} Submitted by AlyciaBHZ on behalf of the Omega
Institute (trureturing project,
\url{https://github.com/the-omega-institute/trureturing}).
All proof code in this submission is newly written; no trureturing code
has been vendored.

\end{document}
